\documentclass{situationist-preprint}

\title{Against Bare Particulars}
\msid{MET-23-a}
\msfamily{MET-23}
\msdiscipline{Metaphysics}
\mstype{Principal manuscript}
\mskeywords{bare particulars; substratum; substance; individuation; essence; bundle theory}

\begin{document}
\maketitle

\begin{abstract}
A bare particular is the constituent of an object that bears its properties while having, in itself, none. Bare particulars have been posited to bear properties, to individuate objects that share all their qualities, and to persist through change. We argue that there are none. An object stripped of all its properties is nothing at all, since there is no way of being a particular that is not a way of having properties; and bare substrata are explanatorily idle, since each role assigned to them is either played by the object itself or requires a primitive distinctness that the object already has. The failure of bundle theories based on absolute or relative discernibility (MET-29-a) does not support substrata, because what an object has beyond its qualities is not a component but the object itself: a substance with a real definition, some of whose properties are essential to it.
\end{abstract}

\section{The question}

Locke described the idea of substance in general as the supposition of something, he knew not what, that supports the qualities we observe \citep[II.xxiii.2]{LockeEssay}. \citet{Bergmann1967} and \citet{Allaire1963} gave the idea its modern form: an ordinary object is a complex of universals and a bare particular, which instantiates the universals and individuates the object, and which has no nature of its own. \citet{Sider2006} defended a modest version on which the particular is \enquote{bare} only in that its properties are not among its parts. The alternatives are the bundle theory, on which an object is nothing over and above its properties, and the substance view, on which the object itself, not a component of it, is what has the properties \citep{Loux1978,Lowe2006}.

\section{Framework}

\begin{definition}\label{d:bare}
A \emph{bare particular} is an entity that instantiates properties and has none of them in virtue of what it is: no property is part of its nature.
\end{definition}

\begin{definition}\label{d:roles}
The \emph{roles} claimed for bare particulars are: (i) to be what bears an object's properties; (ii) to make two objects with the same qualities two; (iii) to be what persists when an object changes its properties.
\end{definition}

\begin{definition}\label{d:substance}
A \emph{substance} is an object that has properties and has a real definition stating what it is; the properties stated in its real definition are essential to it (MET-13).
\end{definition}

\section{The argument}

\begin{proposition}\label{p:nothing}
An object stripped of all its properties is nothing at all.
\end{proposition}

\begin{proof}
Suppose an entity has no property in virtue of what it is. Then nothing about it, considered in itself, distinguishes it from any other entity, from the absence of an entity, or from a different number of entities: there is no answer to what it is, and so no answer to which entity it is or whether it is one. An entity of which there is no answer to what it is lacks the minimum required to be something rather than nothing, since being something is being something or other. A bare particular in the sense of Definition~\ref{d:bare} is therefore not an entity, and the supposed remainder of an object stripped of its properties is nothing.
\end{proof}

The defender of substrata will reply that the bare particular has properties, namely those it instantiates, and is bare only in that none belongs to its nature. That is Sider's version, and it gives up bareness in the sense that matters: the particular is then an entity that has properties, and the question is only whether some of them belong to its nature. If none do, Proposition~\ref{p:nothing} applies to what it is in itself; if some do, it is a substance.

\begin{proposition}\label{p:idle}
Each role of Definition~\ref{d:roles} is played by the object itself, or requires a primitive distinctness that the object already has. Bare particulars are therefore explanatorily idle.
\end{proposition}

\begin{proof}
(i) The bearer of an object's properties is the object: the rose is red, not some constituent of it. Introducing a constituent that bears the redness on the rose's behalf does not explain how anything bears a property; it moves the question from the rose to the constituent. (ii) Two objects with the same qualities, such as Black's spheres, are made two, on the substratum view, by having different substrata. But the substrata differ in no quality either, so their distinctness is primitive. The substance view takes the distinctness of the two spheres as primitive in the same way; the substratum view adds a second entity whose distinctness is equally primitive and does no further work. (iii) What persists when the rose fades is the rose. The substratum view says that the substratum persists and the rose's properties change; but the substratum, having no nature, has no persistence conditions of its own, and the identity of the rose over time is fixed, if by anything, by its kind and its real definition, which belong to the rose and not to the substratum.
\end{proof}

\begin{theorem}\label{t:main}
There are no bare particulars.
\end{theorem}

\begin{proof}
By Proposition~\ref{p:nothing}, a bare particular in the strict sense is nothing; by Proposition~\ref{p:idle}, a substratum in any weaker sense does no work that the object does not already do, and so there is no reason to posit it.
\end{proof}

\section{Objections and replies}

\begin{objection}[The argument from indiscernibles]
MET-29-a shows that bundle theories based on absolute or relative discernibility fail for Black's spheres: two objects can share all their qualities. So objects have something beyond their qualities, and that is the substratum.
\end{objection}

\begin{reply}
The premise is correct and the inference is not. What the spheres have beyond their qualities is not a further constituent but their being two substances. The bundle theory fails because qualities cannot make two things two; a substratum cannot make them two either, except by being two primitively, which the spheres themselves already are. MET-29-a also leaves open a bundle theory based on weak discernibility; the present argument does not depend on rejecting it, since it concerns only what an object would be if its properties were removed. Our own view of objects is the substance view: objects instantiate the sparse universals of MET-02-a, and are neither bundles of them nor universals plus substrata.
\end{reply}

\begin{objection}[Substance is the substratum renamed]
To say that the substance, rather than a substratum, bears the properties is a verbal change. The substance is what remains once its properties are considered apart, and that is the substratum.
\end{objection}

\begin{reply}
The difference is that a substance has a nature. By Definition~\ref{d:substance}, it has a real definition stating what it is, and the properties in that definition are essential to it; so a substance cannot be considered apart from all its properties without ceasing to be considered at all. A substratum, by definition, has no such properties. The substance view thus differs from the substratum view exactly on the point at which Proposition~\ref{p:nothing} applies.
\end{reply}

\section{Conclusion}

An object stripped of all its properties is nothing, and the roles assigned to bare substrata are played by objects themselves or require a primitive distinctness objects already have. There are no bare particulars. Together with MET-02-a, MET-08-a and MET-29-a, the result places the bearers of the sparse universals neither in bundles nor in substrata but in substances with essential properties (MET-13).

\begin{thebibliography}{9}\small
\bibitem[Allaire(1963)]{Allaire1963} Allaire, E.~B. (1963). Bare particulars. \emph{Philosophical Studies}, 14(1--2), 1--8.
\bibitem[Bergmann(1967)]{Bergmann1967} Bergmann, G. (1967). \emph{Realism: A Critique of Brentano and Meinong}. University of Wisconsin Press.
\bibitem[Locke(1690)]{LockeEssay} Locke, J. (1690). \emph{An Essay Concerning Human Understanding}.
\bibitem[Loux(1978)]{Loux1978} Loux, M.~J. (1978). \emph{Substance and Attribute}. Reidel.
\bibitem[Lowe(2006)]{Lowe2006} Lowe, E.~J. (2006). \emph{The Four-Category Ontology}. Oxford University Press.
\bibitem[Sider(2006)]{Sider2006} Sider, T. (2006). Bare particulars. \emph{Philosophical Perspectives}, 20, 387--397.
\end{thebibliography}

\end{document}
